\chapter{Remaining Work and Future Development}
\label{ch:future}

\section{Tasks Remaining Before Final Submission}

Although the core validation pipeline is functional and the experimental results presented
in \cref{ch:results} demonstrate its feasibility, several components require further
development before the system can be considered production-ready.

\subsection{Expanded Dataset and Generalisation Testing}

The current dataset covers 20 template types with 150 samples. Before final submission,
the dataset will be expanded to include all available template variants (estimated
25–30 types) and real failure samples from IKEA's QA team where available. Per-template
type performance will be reported separately to identify template types where the system
underperforms.

\subsection{Colour-Histogram Validation Module}

As identified in the error mode analysis (\cref{sec:experiments}), colour-mismatch
violations are not reliably detected by any of the four current methods. A
colour-histogram comparison module will be added as a fifth validation method in the
ensemble. The module will compute per-zone HSV histogram distances between the generated
label and the reference and flag zones where the histogram divergence (Jensen-Shannon
distance) exceeds a calibrated threshold.

\subsection{Layout Checker for Zone Ordering}

A dedicated zone-ordering validation component will verify that named spatial zones
(product identity, barcode, compliance, date) appear in the expected left-to-right
sequence. This requires zone segmentation, for which the LayoutParser toolkit~\cite{shen2021layoutparser}
with a fine-tuned Detectron2 model will be evaluated.

\subsection{Dashboard Visualisation}

A browser-based dashboard is planned that renders the validation report as a visual
overlay—highlighting violated regions in red on the label image and listing rule
failures in a structured table. Initial prototyping using React/TypeScript is
underway; the final dashboard will be integrated with the FastAPI REST endpoint.

\subsection{Completeness and Cross-Field Consistency Checks}

The metadata validation currently checks field presence and format independently.
Two additional cross-field checks will be implemented:
\begin{itemize}
  \item \textbf{Barcode content vs.\ printed text}: the article number encoded in the
  GS1 DataMatrix must match the human-readable article number printed on the label face.
  \item \textbf{Conditional field presence}: certain fields are required only when
  other fields are present (e.g.\ the YYWW date stamp is required when the date zone
  is activated). These conditional rules will be added to the JSON schema files.
\end{itemize}

\subsection{Full Evaluation on Held-Out Production Labels}

After all components are complete, a final evaluation on a fresh set of production labels
(not used during any development phase) will be conducted. The primary metric is
precision on the non-compliant class, as false positives (flagging a correct label as
invalid) are the most disruptive error mode in an industrial workflow.

\section{Long-Term Future Development}

Beyond the scope of this thesis, several directions are identified for future research
and production development.

\subsection{Fine-Tuned LayoutLM Classifier}

The exported training records (generated label $+$ field annotations $+$ compliance
label) form a dataset suitable for fine-tuning a LayoutLMv3~\cite{huang2022layoutlmv3}
model to perform joint text-layout-image validation. Such a model could capture
semantic content violations (e.g.\ a product name that is present but refers to the
wrong product) that the current visual comparison pipeline cannot detect.

\subsection{PatchCore-Based Anomaly Detection}

The PatchCore method~\cite{roth2022patchcore} stores a coreset of normal patch
embeddings and detects anomalies by nearest-neighbour distance. Adapting PatchCore
to the label validation domain would allow anomaly detection without any negative
training examples—an important property given that faulty labels are not archived.

\subsection{Integration with Product Information Management System}

A production-grade deployment would retrieve the expected field values (article number,
origin, DataMatrix content) directly from IKEA's Product Information Management (PIM)
system via API, enabling content-correctness validation (not just format correctness).

\subsection{Continuous Learning Loop}

As QA operators correct false positives and false negatives in the dashboard, their
corrections could be fed back as additional training examples for the ML classifiers
through the export-JSONL mechanism. This would allow the system to improve over time
without requiring manual dataset curation.
